\chapter{Discusión}

\section{Búsqueda de hiperparámetros}

El HV medio alcanzado por las mejores configuraciones depende principalmente de la dificultad del problema de clasificación y no de la dimensionalidad de entrada (aproximadamente 0.93 en \texttt{ds2}, 0.87 en \texttt{ds1}, 0.81 en \texttt{ds4} y 0.75 en \texttt{ds3}, sin relación monótona con el número de atributos), algo esperable dado que el HV combina simultáneamente exactitud y complejidad estructural. El análisis de importancia de hiperparámetros muestra que la escala de pesos (\texttt{wScale}) es sistemáticamente el hiperparámetro más influyente en tres de los cuatro estudios, coherente con su papel de amplificador necesario para que las conexiones codificadas entre [0,1] generen disparos efectivos en la red; las excepciones son \texttt{ds2}, dominado por la variable \texttt{architecture}, y \texttt{ds4}, donde la importancia se reparte entre varios hiperparámetros sin que ninguno domine, en un espacio de mayor dimensionalidad. Bajo un mismo presupuesto de búsqueda, MOEA-CKF y S-NSGA-II alcanzan HV prácticamente equivalentes en la configuración ganadora: en tres de los cuatro datasets (\texttt{ds1}, \texttt{ds3} y \texttt{ds4}) S-NSGA-II obtiene un HV medio ligeramente superior, con diferencias entre 0.002 y 0.016, mientras que en \texttt{ds2} la relación se invierte levemente a favor de MOEA-CKF, con una diferencia de apenas 0.0003; esto sugiere que, bajo el mismo presupuesto, ambos SMOEA podrían alcanzar HV comparables, aunque esta lectura se basa en el mejor \emph{trial} de un único estudio y se contrasta con mayor rigor estadístico en el Estudio 1. Estos resultados justifican tratar \texttt{architecture}, \texttt{wScale} y los límites de esparsidad como hiperparámetros externos del pipeline y sustentan las configuraciones utilizadas en la fase de comparación de algoritmos.

\section{Comparación de algoritmos}
\subsection{Estudio 1: Comparación de algoritmos en SNN}
\label{sec:discusion-estudio1}

\paragraph{Hipervolumen}

En hipervolumen (HV), S-NSGA-II obtiene una ventaja estadísticamente significativa en dos de los cuatro datasets (\texttt{ds1}, \texttt{ds3}), MOEA-CKF en uno (\texttt{ds4}), y en el restante (\texttt{ds2}) no se observa diferencia significativa. La magnitud de esta ventaja varía sustancialmente entre datasets: es más marcada en el de mayor dimensionalidad, \texttt{ds4} —donde además se invierte a favor de MOEA-CKF—, y es prácticamente inexistente en \texttt{ds2} (18 características), donde ambos algoritmos convergen a una región estrecha y muy alta del espacio de soluciones (diferencia mediana de solo 0.00016 entre medianas, con márgenes mínimos en las corridas donde gana S-NSGA-II y márgenes mayores en las pocas donde gana MOEA-CKF). Estos patrones son consistentes, aunque no concluyentes con solo cuatro puntos de datos, con la hipótesis de que S-NSGA-II rinde mejor en problemas de menor dimensionalidad (14 y 24 características) y MOEA-CKF en el de mayor dimensionalidad (60 características). Adicionalmente, en tres de los cuatro datasets (\texttt{ds1}, \texttt{ds2}, \texttt{ds3}) MOEA-CKF muestra menor varianza entre corridas que S-NSGA-II, incluso cuando pierde en mediana: tiende a ser más consistente y predecible entre ejecuciones, una propiedad relevante si el criterio de evaluación valora la reproducibilidad además del desempeño puntual.

\paragraph{Frentes de Pareto}

El patrón de complejidad-error observado en los frentes explica en gran medida las diferencias de HV descritas arriba: en los datasets donde gana S-NSGA-II (\texttt{ds1}, \texttt{ds3}), sus redes de mejor accuracy son sistemáticamente más densas y más exactas que las de MOEA-CKF, mientras que en \texttt{ds4}, donde gana MOEA-CKF, los roles se invierten y es este último quien produce redes más densas y más precisas. La cobertura mutua de los frentes combinados (C-metric) refuerza esta lectura: es relativamente pareja en \texttt{ds2} y \texttt{ds3} (en torno a 0.50/0.50 y 0.31/0.33) y marcadamente asimétrica en \texttt{ds4} (0.86/0.14 a favor de MOEA-CKF), de modo que el dataset con la ventaja de HV más contundente es también el de mayor dominancia punto a punto.

Sin embargo, HV y dominancia punto a punto no siempre coinciden: en \texttt{ds1}, pese a tener menos puntos en su frente combinado, el frente de MOEA-CKF domina en mayor proporción al de S-NSGA-II que al revés (55.6\% frente a 16.7\%), lo que indica que la ventaja de HV de S-NSGA-II proviene de cubrir una región más amplia del espacio de soluciones —en particular la banda de mayor complejidad, donde concentra la mayoría de sus puntos— y no de que cada solución individual sea mejor que su contraparte. Un mecanismo relacionado aparece en \texttt{ds3}, donde S-NSGA-II traduce el incremento de complejidad estructural en una reducción más sostenida del error a partir de cierto umbral, lo que explica que su frente combinado sea sustancialmente más grande y contribuye a su ventaja de HV. En \texttt{ds2}, en cambio, S-NSGA-II parece especializarse en soluciones muy esparsas con buenos resultados, mientras que MOEA-CKF explora una región más amplia sin obtener una ganancia proporcional, lo cual es coherente con que el HV no distinga a los algoritmos en ese dataset.

Dos limitaciones metodológicas condicionan esta lectura. Primero, el análisis de frentes se apoya en la solución de menor error de entrenamiento de cada corrida como sustituto de la accuracy de test, ya que el 100\% de las corridas archivadas usan arquitecturas de dos salidas (WTA) y no es posible medir la generalización de forma directa; se trata de una limitación estructural, no de un dato pendiente de calcular, y el análisis complejidad-error de todo el frente es el sustituto disponible. Segundo, el C-metric sobre los frentes combinados debe interpretarse como una magnitud aproximada y no como una proporción de alta precisión: tras deduplicar y filtrar soluciones no dominadas, estos frentes quedan reducidos a entre 4 y 14 puntos, lo que limita la resolución de la métrica en conjuntos tan pequeños.

\paragraph{Convergencia}

En los tres datasets donde existe una diferencia clara de HV (\texttt{ds1}, \texttt{ds3}, \texttt{ds4}), el algoritmo con menor HV final converge sistemáticamente más rápido a su propio techo, mientras que el de mayor HV final asciende más lentamente hacia uno más alto: MOEA-CKF llega antes a un techo más bajo en \texttt{ds1} y \texttt{ds3} ($\sim4.3\times$ más rápido en \texttt{ds1}), mientras que en \texttt{ds4} es S-NSGA-II quien se estabiliza antes en un nivel inferior y MOEA-CKF quien tarda $\sim2.4\times$ más en alcanzar el suyo, más alto, sin llegar a estabilizarse del todo dentro del presupuesto de evaluaciones. Un \textit{plateau} rápido en un techo bajo frente a un ascenso lento hacia un techo alto parece ser, entonces, un balance recurrente en este estudio, y no el ruido de una corrida aislada. En \texttt{ds2}, la ausencia de una diferencia relevante de HV se refleja también en la convergencia: la población inicial ya alcanza el 90\% del HV final para ambos algoritmos, por lo que la velocidad de convergencia mide más bien qué tan rápido se estabiliza la búsqueda que una carrera real de optimización sobre un problema comparativamente saturado.

Una limitación transversal a este análisis es que el eje de evaluaciones (FE) de las curvas y umbrales de velocidad es una aproximación por interpolación lineal dentro de cada corrida: no se registra el FE por fila sino solo la \texttt{generation}, que no es directamente comparable entre algoritmos porque MOEA-CKF consume más FE por generación que S-NSGA-II.

\subsection{Estudio 2: Plataforma SNN (C++) vs.\ ANN (MATLAB)}
\label{sec:discusion-estudio2}

Es el hallazgo más uniforme de todo el análisis: la implementación ANN original de PlatEMO/MATLAB supera a la SNN en HV y exactitud en las ocho combinaciones de dataset y algoritmo evaluadas, siempre de forma significativa tras corrección. La brecha de exactitud es menor en \texttt{ds1} y mayor en \texttt{ds4} para S-NSGA-II. La brecha de complejidad muestra, además, un patrón dependiente del algoritmo que se repite entre datasets: para S-NSGA-II, las redes ANN de mejor exactitud tienden a ser casi completamente densas frente a redes SNN mucho más esparsas (particularmente marcado en \texttt{ds2} y \texttt{ds3}), mientras que para MOEA-CKF esta brecha es sistemáticamente más moderada. La única excepción a este patrón ocurre en \texttt{ds4}, donde es la red SNN de MOEA-CKF la que resulta más densa que la ANN, la única anomalía de complejidad entre las ocho combinaciones plataforma-algoritmo del estudio. En conjunto, estos resultados sugieren que la SNN podría favorecer soluciones más esparsas que su contraparte ANN para una parte importante de las combinaciones algoritmo-dataset, un patrón dependiente del algoritmo que podría merecer investigación aparte.

MATLAB/PlatEMO original y la migración a C++ difieren en más que el modelo de neurona, ya que también difieren en la implementación y en el generador de números aleatorios. Por ello, los resultados de esta comparación deben leerse como evidencia de que la implementación ANN concreta supera consistentemente a la implementación SNN concreta evaluada aquí, y no como una conclusión general sobre la neurona de disparo frente a una neurona de tasa.

\subsection{Otros resultados}
\label{sec:discusion-otros}

El efecto del remuestreo sobre el HV final es heterogéneo entre datasets y algoritmos, sin que se identifique una regla simple. Mejora el HV de forma clara y con efecto grande en ambos algoritmos en \texttt{ds1} y \texttt{ds2} —este último con el menor costo relativo de tiempo de los cuatro datasets ($\times4.8$--$4.9$, frente a $\times5$--$15$ en el resto)—, pero produce resultados mixtos en \texttt{ds3} y \texttt{ds4}: en \texttt{ds3} solo mejora a MOEA-CKF, mientras que en S-NSGA-II no hay efecto significativo ($p=0.066$), con la dirección del tamaño de efecto incluso sugiriendo que la evaluación normal iguala o supera al remuestreo más seguido que a la inversa; en \texttt{ds4} el patrón se invierte y es solo S-NSGA-II quien mejora significativamente. El costo de tiempo del remuestreo, en cambio, es siempre sustancial (entre $\times4.8$ y $\times15.1$) independientemente de si el HV mejora o no, lo que deja como pregunta abierta si existe alguna interacción entre el remuestreo y los mecanismos específicos de cada algoritmo que explique esta heterogeneidad.

En tiempo de cómputo, S-NSGA-II es significativamente más rápido que MOEA-CKF en los cuatro datasets, con efecto grande y menor variabilidad entre corridas en todos los casos. La brecha, sin embargo, crece marcadamente con la dimensionalidad del problema: de aproximadamente $6$--$10\%$ en \texttt{ds1}--\texttt{ds3} (14--24 características) a cerca de $2.4\times$ en \texttt{ds4} (60 características, sin una sola excepción en las 15 corridas pareadas), donde MOEA-CKF resulta además el más lento y variable. Este patrón es consistente con que el mecanismo de fusión de conocimiento entre escalas de MOEA-CKF (SVD, análisis de dispersión) consume más FE por generación y escala peor con la dimensionalidad que los operadores más simples de S-NSGA-II.

Estas comparaciones de tiempo solo son válidas dentro del conjunto experimental temporal, ejecutado con hiperparámetros compartidos entre ambos algoritmos; los conjuntos experimentales normal y de remuestreo no permiten una comparación válida en tiempo de ejecución debido al paralelismo y los hilos no controlados entre tandas de ejecución, por lo que sus conclusiones de tiempo no deben mezclarse con las conclusiones de calidad (HV) del resto del análisis, que sí emplea hiperparámetros afinados.

\section{Conectividad de las redes}
\label{sec:discusion-conectividad-v2}

Los dos indicadores de la Sección~\ref{sec:conectividad} —clases de salida desconectadas y conexiones inútiles— responden a mecanismos distintos y conviene discutirlos por separado, aunque ambos comparten una misma raíz: ninguno de los operadores de dispersión de MOEACKF ni de S-NSGA-II tiene noción alguna de a qué neurona o a qué salida pertenece cada sinapsis, y ninguno de los dos objetivos de la optimización ($f_1$, complejidad; $f_2$, error de entrenamiento) penaliza directamente una sinapsis activa que no forma parte de ningún camino completo entre la entrada y la salida.

\subsection{Clases de salida desconectadas: desbalance de clases}

La asimetría por clase ya señalada en el análisis de resultados —en \texttt{ds2} es siempre la clase 0 la que pierde su conexión, en \texttt{ds3} es siempre la clase 1— coincide exactamente con la clase minoritaria de cada dataset: en \texttt{ds2} (Climate) la clase 0 representa apenas 46 de 540 instancias (8.5\%); en \texttt{ds3} (German) la clase minoritaria representa 300 de 1000 (30\%). En \texttt{ds1} (Australian, clase 0/clase 1: 55.5\%/44.5\%) y \texttt{ds4} (Sonar, 53.4\%/46.6\%), donde ambas clases están razonablemente balanceadas, el porcentaje de redes con una salida desconectada es nulo o casi nulo (0\% en tres de las cuatro combinaciones, 6.5\% en la restante). Esto es consistente con que, al no distinguir los operadores de dispersión entre clases al apagar pesos, basta con que $f_2$ (el error de entrenamiento) recompense predecir siempre la clase mayoritaria —una estrategia que en \texttt{ds2} ya alcanza 91.5\% de exactitud sin necesitar ninguna conexión hacia la clase minoritaria— para que la búsqueda converja hacia soluciones que sacrifican por completo la conectividad de esa clase, sin que el objetivo de exactitud lo impida.

La excepción parcial es \texttt{ds3}$\times$SNSGAII, con 0\% de salidas desconectadas pese al mismo desbalance 70/30 que en \texttt{ds3}$\times$MOEACKF produce 80.6\%. La diferencia es coherente con el diseño de cada operador: el operador de mutación de dispersión de S-NSGA-II muta un nivel de dispersión objetivo por individuo y apaga o reactiva pesos al azar hasta alcanzarlo, mientras que MOEACKF construye su máscara inicial activando un número aleatorio $k \in [1, D]$ de variables mediante selección por torneo sesgada por aptitud, sin ninguna garantía de que ese conjunto —que puede ser tan pequeño como una sola sinapsis de las $D$ totales— toque la neurona de salida de la clase minoritaria. Con un espacio de decisión de cientos o miles de variables, es más fácil que un sorteo casi puntual dependa del azar de la clase minoritaria que un ajuste incremental de densidad la deje completamente fuera, aunque esta lectura es una asociación entre el mecanismo del operador y el patrón observado, no una prueba directa sobre el código en ejecución.

\subsection{Conexiones inútiles: asimetría estructural y presión de poda débil}

El porcentaje de conexiones inútiles no aparece asociado con el desbalance de clases sino con la dimensionalidad de entrada: promediando ambos algoritmos por dataset (y excluyendo \texttt{ds2}$\times$SNSGAII por la razón que se explica más abajo), el porcentaje sube de 31.5\% en \texttt{ds1} (14 características) a 54.3\% en \texttt{ds3} (24 características) y 57.3\% en \texttt{ds4} (60 características). Esta tendencia es consistente con una explicación estructural propia de la codificación: cada neurona oculta tiene $n_F+1$ posiciones de entrada posibles (una por característica más el sesgo) pero solo $n_O=2$ posiciones de salida posibles, de modo que, ante cualquier mecanismo de activación de pesos sin conciencia de la estructura del grafo, la probabilidad de que una neurona reciba al menos una entrada activa crece con $n_F$ mucho más rápido que la probabilidad de que además tenga una salida activa —fija en solo dos posiciones—. Cuantas más características tiene el dataset, más fácil es entonces que una neurona termine con entrada real pero sin salida (sumidero) que con ambas (viva), lo que aumenta la fracción de sinapsis activas atrapadas en caminos incompletos. Esta lectura es razonable con solo cuatro datasets y ocho combinaciones, y no se contrastó de forma independiente del nivel de dispersión de cada solución.

\texttt{ds2}$\times$SNSGAII es la excepción que confirma esa lectura por la vía opuesta: su solución de mejor exactitud alcanza 92.1\% de accuracy con solo el 0.32\% de las 840 sinapsis activas y una única neurona oculta viva. Con tan pocas sinapsis activas en total, no queda margen combinatorio para que aparezcan neuronas parcialmente conectadas, y el 0\% de conexiones inútiles es entonces un efecto secundario de la extrema dispersión y no evidencia de una red mejor estructurada que el resto.

A esta asimetría estructural se suma una segunda causa, de naturaleza algorítmica: ninguno de los dos objetivos premia explícitamente eliminar una conexión inútil. Retirar una sinapsis que llega a una neurona sumidero (o que sale de una neurona fuente) reduce $f_1$ sin alterar $f_2$ en absoluto, de modo que la solución resultante domina estrictamente a la original y en teoría debería reemplazarla en el frente. En la práctica, esa poda solo ocurre si algún operador de variación llega a generar exactamente esa vecina —con esa sinapsis puntual desactivada y ninguna otra— dentro del presupuesto de evaluaciones disponible, algo cada vez menos probable a medida que $D$ crece (de 680 en \texttt{ds1} a 2520 en \texttt{ds4}) y que ninguno de los operadores dirige la búsqueda hacia esa vecina en particular. La diferencia entre algoritmos observada en la Tabla~\ref{tab:conectividad-algo} (49.5\% en MOEACKF frente a 32.7\% en S-NSGA-II) es consistente con esta lectura: la máscara de MOEACKF activa posiciones sueltas del vector de decisión sin agrupación alguna, mientras que la máscara de S-NSGA-II opera en franjas contiguas que —dado que la codificación agrupa las entradas de cada neurona oculta en un bloque contiguo del vector— tienden a activar o desactivar el bloque completo de una neurona en conjunto, produciendo con algo más de frecuencia neuronas enteramente vivas o enteramente muertas en vez de parcialmente conectadas.

\subsection{Implicancia práctica}

Ambos fenómenos afectan directamente a la red que se desplegaría en la práctica: en la Tabla~\ref{tab:conectividad-algo}, poco más de un tercio de las soluciones de mejor exactitud de MOEACKF y poco más de un cuarto de las de S-NSGA-II son incapaces de predecir una de las dos clases sin importar la entrada, y en promedio entre un tercio y la mitad de las sinapsis activas de esas soluciones no contribuye a ninguna predicción. Ninguno de los dos indicadores se corrige evaluando de mejor manera la red resultante, sino interviniendo en los operadores que deciden qué sinapsis permanecen activas: una inicialización o un operador de reparación que garanticen al menos un camino activo hacia cada clase de salida atacarían directamente el primer problema, y una poda estructural —que retire sinapsis pertenecientes a neuronas sumidero o fuente sin costo en $f_2$— atacaría el segundo; ninguna de las dos intervenciones fue implementada ni evaluada en el marco de esta tesis.
